
\subsubsection{Momentum dynamics}

\paragraph{Global internal boundaries}

Momentum dynamics is described in spherical coordinates, because collisions
is the dominant physical process for current drive and matrices are
therefore expected to be well conditionned. Consequently, internal
boundaries must be specified at $p=0$ and $\left| \xi _{0}\right| =1$ which
corresponds to the conditions 
\begin{equation}
\left\{ 
\begin{array}{l}
f_{0}^{\left( 0\right) }\left( \psi ,-p,\arccos \xi _{0}\right)
=f_{0}^{\left( 0\right) }\left( \psi ,-p,\pi -\arccos \xi _{0}\right) \\ 
f_{0}^{\left( 0\right) }\left( \psi ,p,-\arccos \xi _{0}\right)
=f_{0}^{\left( 0\right) }\left( \psi ,p,\arccos \xi _{0}\right) \\ 
f_{0}^{\left( 0\right) }\left( \psi ,p,\pi +\arccos \xi _{0}\right)
=f_{0}^{\left( 0\right) }\left( \psi ,p,\pi -\arccos \xi _{0}\right)
\end{array}
\right\}
\end{equation}

Since the type of equation to be solved correspond to a Neuman problem,
gradients must be specified at the internal boundaries. To fullfil above
conditions, one must thus determine $\left( \Delta p,\Delta \xi _{0}\right) $
at $p=0$ and $\left| \xi _{0}\right| =1$ for the flux grid only, since the
distribution grid is interlaced by definition with the flux one. Therefore,
starting from grids definition given in Sec. ***, 
\begin{equation}
\Delta \xi _{0,0}=\frac{\Delta \xi _{0,1/2}+\Delta \xi _{0,-1/2}}{2}=\Delta
\xi _{0,1/2}=\xi _{0,1}-\xi _{0,0}=2\xi _{0,1/2}+1
\end{equation}
as $\Delta \xi _{0,1/2}=\Delta \xi _{0,-1/2},$ $\xi _{0,1/2}=\frac{\xi
_{0,1}+\xi _{0,0}}{2}$et $\xi _{0,0}=-1.$

Much in the same way, 
\begin{equation}
\Delta \xi _{0,n_{\xi }}=\frac{\Delta \xi _{0,n_{\xi }+1/2}+\Delta \xi
_{0,n_{\xi }-1/2}}{2}=\Delta \xi _{0,n_{\xi }-1/2}=\xi _{0,n_{\xi }}-\xi
_{0,n_{\xi }-1}=1-\xi _{0,n_{\xi }-1}
\end{equation}
since $\Delta \xi _{0,n_{\xi }+1/2}=\Delta \xi _{0,n_{\xi }-1/2}.$
Furthermore, since by definition, 
\begin{equation}
\xi _{0,n_{\xi }-1/2}=\frac{\xi _{0,n_{\xi }}+\xi _{0,n_{\xi }-1}}{2}=\frac{%
1+\xi _{0,n_{\xi }-1}}{2}
\end{equation}
one obtains finally

\begin{equation}
\Delta \xi _{0,n_{\xi }}=2\left( 1-\xi _{0,n_{\xi }-1/2}\right)
\end{equation}

A similar technique may be used for the momentum grid $p$. Hence 
\begin{equation}
\Delta p_{0}=\frac{1}{2}\left( \Delta p_{1/2}+\Delta p_{-1/2}\right) =\Delta
p_{1/2}
\end{equation}
using $\Delta p_{-1/2}=\Delta p_{1/2}.$ Since $\Delta p_{1/2}=p_{1}-p_{0}$,
and $p_{1/2}=\frac{p_{1}+p_{0}}{2}$ one obtains finally, 
\begin{equation}
\Delta p_{0}=2p_{1/2}
\end{equation}
because $p_{0}=0.$

It is important to note that internal boundary conditions are only needed in
the evaluation of the cross-derivative terms, since in the discrete form of
the Fokker-Planck equation using two grids, as shown, in Sec. ***, it is
automatically fullfiled. Indeed, at fixed $\psi _{l+1/2},$%
\begin{eqnarray}
\left. \frac{\widetilde{q}\left( \psi \right) }{B_{\min }\left( \psi \right) 
}p^{2}\frac{\partial f^{\left( 0\right) }}{\partial p}\right|
_{l+1/2,1/2,j+1/2}^{(k+1)} &=&\frac{\widetilde{q}\left( \psi _{l+1/2}\right) 
}{B_{\min }\left( \psi _{l+1/2}\right) }\left. \frac{\partial \left( p^{2}%
\mathbf{S}_{p}^{\left( 0\right) }\right) }{\partial p}\right|
_{l+1/2,1/2,j+1/2}^{(k+1)}  \nonumber \\
&&-\frac{\widetilde{q}\left( \psi _{l+1/2}\right) }{B_{\min }\left( \psi
_{l+1/2}\right) }\frac{p_{1/2}}{\lambda \left( \psi _{l+1/2},\xi
_{0,j+1/2}\right) }\times  \nonumber \\
&&\left. \frac{\partial }{\partial \xi _{0}}\left( \sqrt{1-\xi _{0}^{2}}%
\lambda \left( \psi ,\xi _{0}\right) \mathbf{S}_{\xi }^{\left( 0\right)
}\right) \right| _{l+1/2,1/2,j+1/2}^{(k+1)}
\end{eqnarray}
which gives

\begin{eqnarray}
\left. \frac{\widetilde{q}\left( \psi \right) }{B_{\min }\left( \psi \right) 
}p^{2}\frac{\partial f^{\left( 0\right) }}{\partial p}\right|
_{l+1/2,1/2,j+1/2}^{(k+1)} &=&\frac{\widetilde{q}\left( \psi _{l+1/2}\right) 
}{B_{\min }\left( \psi _{l+1/2}\right) }\frac{p_{1}^{2}\mathbf{S}%
_{p,l+1/2,1,j+1/2}^{\left( 0\right) (k+1)}}{\Delta p_{1/2}}  \nonumber \\
&&-\frac{\widetilde{q}\left( \psi _{l+1/2}\right) }{B_{\min }\left( \psi
_{l+1/2}\right) }\frac{p_{1/2}}{\lambda \left( \psi _{l+1/2},\xi
_{0,j+1/2}\right) }\times  \nonumber \\
&&\left. \frac{\partial }{\partial \xi _{0}}\left( \sqrt{1-\xi _{0}^{2}}%
\lambda \left( \psi ,\xi _{0}\right) \mathbf{S}_{\xi }^{\left( 0\right)
}\right) \right| _{l+1/2,1/2,j+1/2}^{(k+1)}  \label{boundary_p}
\end{eqnarray}
since $p_{0}^{2}\mathbf{S}_{p,l+1/2,0,j+1/2}^{\left( 0\right) (k+1)}=0$ with
the flux grid here considered, with $p_{0}=0$. Much in the same way, 
\begin{eqnarray}
\left. \frac{\widetilde{q}\left( \psi \right) }{B_{\min }\left( \psi \right) 
}p^{2}\frac{\partial f^{\left( 0\right) }}{\partial p}\right|
_{l+1/2,i+1/2,1/2}^{(k+1)} &=&\frac{\widetilde{q}\left( \psi _{l+1/2}\right) 
}{B_{\min }\left( \psi _{l+1/2}\right) }\left. \frac{\partial \left( p^{2}%
\mathbf{S}_{p}^{\left( 0\right) }\right) }{\partial p}\right|
_{l+1/2,i+1/2,1/2}^{(k+1)}  \nonumber \\
&&-\frac{\widetilde{q}\left( \psi _{l+1/2}\right) }{B_{\min }\left( \psi
_{l+1/2}\right) }\frac{p_{i+1/2}}{\lambda \left( \psi _{l+1/2},\xi
_{0,1/2}\right) }\times  \nonumber \\
&&\frac{\sqrt{1-\xi _{0,1}^{2}}\lambda \left( \psi ,\xi _{0,1}\right) 
\mathbf{S}_{\xi ,l+1/2,i+1/2,1}^{\left( 0\right) (k+1)}}{\Delta \xi _{0,1/2}}
\end{eqnarray}
and 
\begin{eqnarray}
\left. \frac{\widetilde{q}\left( \psi \right) }{B_{\min }\left( \psi \right) 
}p^{2}\frac{\partial f^{\left( 0\right) }}{\partial p}\right|
_{l+1/2,i+1/2,n_{\xi _{0}}-1/2}^{(k+1)} &=&\frac{\widetilde{q}\left( \psi
_{l+1/2}\right) }{B_{\min }\left( \psi _{l+1/2}\right) }\left. \frac{%
\partial \left( p^{2}\mathbf{S}_{p}^{\left( 0\right) }\right) }{\partial p}%
\right| _{l+1/2,i+1/2,n_{\xi _{0}}-1/2}^{(k+1)}  \nonumber \\
&&+\frac{\widetilde{q}\left( \psi _{l+1/2}\right) }{B_{\min }\left( \psi
_{l+1/2}\right) }\frac{p_{i+1/2}}{\lambda \left( \psi _{l+1/2},\xi
_{0,n_{\xi _{0}}-1/2}\right) }\times  \nonumber \\
&&\frac{\sqrt{1-\xi _{0,n_{\xi _{0}}-1}^{2}}\lambda \left( \psi ,\xi
_{0,n_{\xi _{0}}-1}\right) \mathbf{S}_{\xi ,l+1/2,i+1/2,n_{\xi
_{0}}-1}^{\left( 0\right) (k+1)}}{\Delta \xi _{0,n_{\xi _{0}}-1/2}}
\end{eqnarray}
since $\sqrt{1-\xi _{0,0}^{2}}\lambda \left( \psi ,\xi _{0,0}\right) \mathbf{%
S}_{\xi ,l+1/2,i+1/2,0}^{\left( 0\right) }=0$ and $\sqrt{1-\xi _{0,n_{\xi
_{0}}}^{2}}\lambda \left( \psi ,\xi _{0,n_{\xi _{0}}}\right) \mathbf{S}_{\xi
,l+1/2,i+1/2,n_{\xi _{0}}}^{\left( 0\right) }=0,$ because $\xi
_{0,0}^{2}=\xi _{0,n_{\xi _{0}}}^{2}=1$ by construction.

\paragraph{Trapped-passing internal boundaries}

Since the bounce time is much less than all other characteristics times of
physical processes that are involved in the kinetic calculations,

\paragraph{External boundaries}

The other boundaries are inserted into the problem in violation of the true
physical picture, mainly because the momentum domain extends off to
infinity, while only a subspace is considered in computations. The upper
limit of the domain is therefore chosen so that the interesting physics may
be accurately described, i.e. $p_{\max }\gg p_{th}$ for studying the
Maxwellian distribution function, $p_{\max }\gg \gamma \left( \omega
/k\right) _{\max }$ for the RF waves and $p_{\max }\gg p_{Dreicer}$ for the
runaway electrons. These conditions must ensure the conservative nature of
the problem here addressed, i.e. that plasma cannot enter or leave the
domain of integration. This corresponds to the local condition 
\begin{equation}
\mathbf{S}_{\mathbf{p}}\cdot \widehat{n}=0  \label{Externboundary}
\end{equation}
where $\widehat{n}$ is the normal to the external boundary. Since for
symmetry reasons related to the nature of the collision operator, the
subspace considered for computations is a sphere of radius $p_{\max },$ $%
\widehat{n}=\widehat{p}.$ When the condition (\ref{Externboundary}) is
fullfiled, external boundary conditions have almost a negligible influence,
and the numerical solution in the subspace is usually close to the
theoretical one. It is usually the case for the RF problem, since the range
of resonant interaction between the wave and the electrons is well localized
in momentum space, far enough from the boundary so that no electron can
leave the domain of integration.

The Fokker-Planck equation reduces an hyperbolic equation in the vicinity of 
$p_{\max }$ (high velocity limit), and therefore no boundary condition must
be specified, and the standard upstream differencing applies. Indeed, since
for collisions, $D_{pp}^{C\left( 0\right) }/v\sim 1/v^{4},$ while $%
F_{p}^{C\left( 0\right) }\sim 1/v^{2},$ and one can consider that $%
D_{pp}^{C\left( 0\right) }=0$ at $p=p_{\max }.$ Here, the pitch-angle
scattering term $D_{\xi \xi }^{C\left( 0\right) }$requires no special
handling, because it causes parallel diffusion to the boundary.

Nevertheless, since cross-derivative terms potentially involves points
outside the computation domain, the following rule is used, even if the
exact value of $\Delta p_{\max }$ has no importance,

\begin{equation}
\Delta p_{n_{p}}=\Delta p_{\max }=\frac{1}{2}\left( \Delta
p_{n_{p}+1/2}+\Delta p_{n_{p}-1/2}\right) =\Delta
p_{n_{p}-1/2}=p_{n_{p}}-p_{n_{p}-1}
\end{equation}

As by definition $p_{n_{p}+1/2}=\frac{p_{n_{p}}+p_{n_{p}-1}}{2}=\frac{%
p_{\max }+p_{n_{p}-1}}{2},$ it results 
\begin{equation}
\Delta p_{n_{p}}=2\left( p_{n_{p}}-p_{n_{p}-1/2}\right)
\end{equation}
or in equivalent manner, 
\begin{equation}
\Delta p_{\max }=2\left( p_{\max }-p_{n_{p}-1/2}\right)
\end{equation}

If an electric field is present, then in the real problem, some electrons
will runaway and leave the domain of integration. In that case, the
condition $\mathbf{S}_{\mathbf{p}}\cdot \widehat{p}=0$ is no more fullfiled,
leading to an effective loss of electrons. Since the Fokker-planck equation
is still of hyperbolic type, no specific modification must be applied to the
limit of the integration domain. However, one must ensure that the total
number of electrons is kept constant. Several techniques may be used to
avoid a decay of the number of electrons at the runaway rate $\Gamma _{R}$,
as defined in Sec. ***.

A possibility is to substitute $f_{0}^{\left( 0\right) }$ by $f_{0}^{\left(
0\right) }\exp \left[ -\int_{0}^{t}\Gamma _{R}^{\left( 0\right) }\left(
t^{\prime }\right) dt^{\prime }\right] $ in the Fokker-Planck equation, so
that $f_{0}^{\left( 0\right) }\left( t\rightarrow \infty \right) $ is
independent of $t$. In that case, $\overline{\overline{M}}%
_{p,l+1/2,i+1/2,j+1/2}^{\left( 0\right) }$ must just be replaced by $%
\overline{\overline{M}}_{p,l+1/2,i+1/2,j+1/2}^{\left( 0\right) }-\Gamma
_{R,l+1/2}^{\left( 0\right) },$ where $\Gamma _{R}^{\left( 0\right) }$ is
the runaway rate calculated at the minimum of the magnetic field $B$, as
defined in Sec. ***. The Fokker-Planck equation becomes therefore slightly
non-linear, since $\Gamma _{R}^{\left( 0\right) }$ is an integral over the
distribution function $f_{0}^{\left( 0\right) }.$ However, since the
population that is concerned is usually very small, the non-linearity is
fairly weak. Nevertheless, the main drawback of this approach is that it
makes the Fokker-Planck equation non-conservative. Furthermore, it
necessitates to recalculate periodically the matrix to be inverted, that can
be very time consuming. Therefore this method is not considered in the
calculations.

Another technique is to inject cold electrons at $p=0$, in order to
compensate runaway losses and maintain the density. Such an approach has the
main advantage to have a clear meaning. By definition, the number of
electrons leaving the integration domain is simply $\Gamma _{R}^{\left(
0\right) }.$ In principle, without particle losses, $%
p_{0}^{2}S_{p,l+1/2,0,j+1/2}^{\left( 0\right) }=0,$ though $%
S_{p,l+1/2,0,j+1/2}^{\left( 0\right) }$ is infinite at $p=0$. Compensation
of hot electron losses by cold ones leads to the relation 
\begin{eqnarray}
2\pi \sum_{j=0}^{n_{\xi _{0}}-1}p_{0}^{2}S_{p,l+1/2,0,j+1/2}^{\left(
0\right) }\Delta \xi _{0,j+1/2} &=&\Gamma _{R,l+1/2}^{\left( 0\right) } 
\nonumber \\
&=&2\pi \sum_{j=0}^{n_{\xi
_{0}}-1}p_{n_{p}-1/2}^{2}S_{p,l+1/2,n_{p}-1/2,j+1/2}^{\left( 0\right)
}\Delta \xi _{0,j+1/2}
\end{eqnarray}

Since the pitch-angle grid is non-uniform, one must consider $S_{p}^{\left(
0\right) }$ for each pitch-angle value, and one obtains 
\begin{equation}
p_{0}^{2}S_{p,l+1/2,0,j+1/2}^{\left( 0\right)
}=p_{n_{p}-1/2}^{2}S_{p,l+1/2,n_{p}-1/2,j+1/2}^{\left( 0\right) }
\end{equation}

Therefore, expression (\ref{boundary_p}) is modified according to the
relation 
\begin{eqnarray}
\left. \frac{\widetilde{q}\left( \psi \right) }{B_{\min }\left( \psi \right) 
}p^{2}\frac{\partial f^{\left( 0\right) }}{\partial p}\right|
_{l+1/2,1/2,j+1/2}^{(k+1)} &=&\frac{\widetilde{q}\left( \psi _{l+1/2}\right) 
}{B_{\min }\left( \psi _{l+1/2}\right) }\left( \frac{p_{1}^{2}\mathbf{S}%
_{p,l+1/2,1,j+1/2}^{\left( 0\right) (k+1)}}{\Delta p_{1/2}}-\frac{p_{0}^{2}%
\mathbf{S}_{p,l+1/2,0,j+1/2}^{\left( 0\right) (k+1)}}{\Delta p_{1/2}}\right)
\nonumber \\
&&-\frac{\widetilde{q}\left( \psi _{l+1/2}\right) }{B_{\min }\left( \psi
_{l+1/2}\right) }\frac{p_{1/2}}{\lambda \left( \psi _{l+1/2},\xi
_{0,j+1/2}\right) }\times  \nonumber \\
&&\left. \frac{\partial }{\partial \xi _{0}}\left( \sqrt{1-\xi _{0}^{2}}%
\lambda \left( \psi ,\xi _{0}\right) \mathbf{S}_{\xi }^{\left( 0\right)
}\right) \right| _{l+1/2,1/2,j+1/2}^{(k+1)}
\end{eqnarray}
or 
\begin{eqnarray}
\left. \frac{\widetilde{q}\left( \psi \right) }{B_{\min }\left( \psi \right) 
}p^{2}\frac{\partial f^{\left( 0\right) }}{\partial p}\right|
_{l+1/2,1/2,j+1/2}^{(k+1)} &=&\frac{\widetilde{q}\left( \psi _{l+1/2}\right) 
}{B_{\min }\left( \psi _{l+1/2}\right) }\left( \frac{p_{1}^{2}\mathbf{S}%
_{p,l+1/2,1,j+1/2}^{\left( 0\right) (k+1)}}{\Delta p_{1/2}}-\frac{%
p_{n_{p}-1/2}^{2}S_{p,l+1/2,n_{p}-1/2,j+1/2}^{\left( 0\right) (k+1)}}{\Delta
p_{1/2}}\right)  \nonumber \\
&&-\frac{\widetilde{q}\left( \psi _{l+1/2}\right) }{B_{\min }\left( \psi
_{l+1/2}\right) }\frac{p_{1/2}}{\lambda \left( \psi _{l+1/2},\xi
_{0,j+1/2}\right) }\times  \nonumber \\
&&\left. \frac{\partial }{\partial \xi _{0}}\left( \sqrt{1-\xi _{0}^{2}}%
\lambda \left( \psi ,\xi _{0}\right) \mathbf{S}_{\xi }^{\left( 0\right)
}\right) \right| _{l+1/2,1/2,j+1/2}^{(k+1)}
\end{eqnarray}

Since $\Gamma _{R,l+1/2}^{\left( 0\right) (k+1)}$ is a weak function of $%
f_{0}^{\left( 0\right) },$ it is possible to replace $%
S_{p,l+1/2,n_{p}-1/2,j+1/2}^{\left( 0\right) (k+1)}$ by its explicit form $%
S_{p,l+1/2,n_{p}-1/2,j+1/2}^{\left( 0\right) (k)}$, so that one may avoid to
recalculate the matrix $\overline{\overline{M}}_{p,l+1/2,i+1/2,j+1/2}^{%
\left( 0\right) }$ at each iteration which is extremely computer time
consuming. Consequently, with this scheme, the zero order Fokker-Planck
equation becomes 
\begin{eqnarray}
&&\frac{\widetilde{q}\left( \psi _{l+1/2}\right) }{B_{\min }\left( \psi
_{l+1/2}\right) }p_{i+1/2}^{2}\frac{f_{0,l+1/2,i+1/2,j+1/2}^{\left( 0\right)
(k+1)}-f_{0,l+1/2,i+1/2,j+1/2}^{\left( 0\right) (k)}}{\Delta t}  \nonumber \\
&&+\left. \frac{\widetilde{q}\left( \psi \right) }{B_{\min }\left( \psi
\right) }p^{2}\frac{\partial f^{\left( 0\right) }}{\partial p}\right|
_{l+1/2,i+1/2,j+1/2}^{(k+1)}  \nonumber \\
&&+\left. \frac{\widetilde{q}\left( \psi \right) }{B_{\min }\left( \psi
\right) }p^{2}\frac{\partial f^{\left( 0\right) }}{\partial \psi }\right|
_{l+1/2,i+1/2,j+1/2}^{(k+1)}  \nonumber \\
&&-\frac{\widetilde{q}\left( \psi _{l+1/2}\right) }{B_{\min }\left( \psi
_{l+1/2}\right) }\frac{p_{n_{p}-1/2}^{2}S_{p,l+1/2,n_{p}-1/2,j+1/2}^{\left(
0\right) (k)}}{\Delta p_{1/2}}  \nonumber \\
&=&0
\end{eqnarray}

\subsubsection{Spatial dynamics}

For the spatial dynamics, 
